\chapter{تحلیل داده و یافته‌های اکتشافی}
\label{ch:data_analysis}

این فصل تصویری منسجم از کوهورت \lr{MIMIC\text{-}IV} ارائه می‌کند: توزیع داده‌ها، همبستگی ویژگی‌ها، کیفیت داده و چند تصویرسازی اکتشافی که مبنای تصمیم‌های مدل‌سازی در فصل‌های بعد هستند.

%==============================================================================
\section{طرحواره و گردش دادۀ \lr{MIMIC\text{-}IV}}
%==============================================================================

شکل~\ref{fig:mimic_schema} رابطۀ جدول‌های \lr{MIMIC\text{-}IV} و روند پیش‌پردازش را به‌صورت خلاصه نشان می‌دهد.

\begin{figure}[htbp]
\centering
\begin{latin}
% MIMIC-IV extraction schema — monochrome, publication-style (TikZ)
% Preamble (once per document): \usepackage{tikz}\usetikzlibrary{positioning,arrows.meta,calc}
\begin{tikzpicture}[
  scale=0.88,
  transform shape,
  font=\footnotesize,
  >=Stealth,
  tbl/.style={
    rectangle,
    draw=black!55,
    line width=0.4pt,
    fill=black!3,
    minimum height=8mm,
    inner xsep=4pt,
    inner ysep=3pt,
    align=center,
    rounded corners=1.5pt,
  },
  code/.style={tbl, font=\footnotesize\ttfamily},
  proc/.style={
    rectangle,
    draw=black!55,
    line width=0.45pt,
    fill=black!6,
    inner sep=6pt,
    align=center,
    rounded corners=2pt,
    font=\footnotesize,
  },
  outbox/.style={
    rectangle,
    draw=black!55,
    line width=0.45pt,
    fill=black!10,
    inner sep=6pt,
    align=center,
    rounded corners=2pt,
    font=\footnotesize\bfseries,
  },
  ed/.style={black!50, line width=0.35pt},
  el/.style={font=\footnotesize, text=black!65, inner sep=1pt, fill=white},
]

% --- Core hospital tables ---
\node[code] (pat) {patients};
\node[code, right=9mm of pat] (adm) {admissions};
\node[code, right=9mm of adm] (icu) {icustays};
\node[code, above=9mm of icu, xshift=-6mm] (dx) {diagnoses\_icd};

\draw[ed, ->] (pat) -- node[el, above=0.5pt] {\texttt{subject\_id}} (adm);
\draw[ed, ->] (adm) -- node[el, above=0.5pt] {\texttt{hadm\_id}} (icu);
\draw[ed, ->] (adm.north) |- node[el, pos=0.28, left=1pt] {\texttt{hadm\_id}} (dx.west);

% --- Labs / chart (second row) ---
\node[code, below left=12mm and 2mm of icu] (dlab) {d\_labitems};
\node[code, right=5mm of dlab] (lab) {labevents};
\node[code, right=7mm of lab] (chart) {chartevents};
\node[code, right=5mm of chart] (dit) {d\_items};

\draw[ed, ->] (dlab) -- node[el, above=0.5pt] {\texttt{itemid}} (lab);
\draw[ed, ->] (dit) -- node[el, above=0.5pt] {\texttt{itemid}} (chart);

\coordinate (hub) at ($(icu.south)+(0,-2.5mm)$);
\draw[ed, ->] (icu.south) -- (hub);
\draw[ed, ->] (hub) -| node[el, pos=0.18, below=1pt] {\texttt{stay\_id}} (lab.north);
\draw[ed, ->] (hub) -| (chart.north);

% --- Preprocessing ---
\node[proc, below=11mm of lab, anchor=north, xshift=16mm] (prep) {%
  \textbf{Preprocessing pipeline}\\[0.35em]
  {\footnotesize 15 labs + 7 vitals $\rightarrow$ clip / min--max $[0,1]$ $\rightarrow$ feature engineering}\\[0.2em]
  {\footnotesize $\rightarrow$ \textbf{25-D tabular} $\;|\;$ \textbf{token sequence} $\rightarrow$ ETHOS $\rightarrow$ \textbf{128-D}}%
};

\coordinate (merge) at ($(prep.north)+(0,4.5mm)$);
\draw[ed, ->] (lab.south) |- (merge);
\draw[ed, ->] (chart.south) |- (merge);
\draw[ed, ->] (merge) -- (prep.north);

% --- Cohort output ---
\node[outbox, below=6mm of prep] (cohort) {$2{,}000$ patients $\times$ $25$ features + labels};
\draw[ed, ->] (prep) -- (cohort);

\end{tikzpicture}

\end{latin}
\caption{طرح استخراج \lr{MIMIC\text{-}IV}: جدول‌های منبع، کلیدهای اتصال، و پیش‌پردازش به بازنمایی‌های جدولی و توکنی (نمودار برداری؛ تک‌رنگ برای چاپ).}
\label{fig:mimic_schema}
\end{figure}

جدول~\ref{tab:mimic_flow} روند استخراج و پیش‌پردازش دادۀ \lr{MIMIC\text{-}IV} را خلاصه می‌کند. ابتدا \texttt{\lr{patients}}، \texttt{\lr{admissions}} و \texttt{\lr{icustays}} بر اساس \texttt{\lr{subject\_id}} و \texttt{\lr{hadm\_id}} به هم متصل می‌شوند؛ سپس \texttt{\lr{labevents}} و \texttt{\lr{chartevents}} از طریق \texttt{\lr{stay\_id}} به اقامت‌های \lr{ICU} پیوند می‌خورند. در مرحلۀ پیش‌پردازش، \pnum{15} آزمایش و \pnum{7} علامت حیاتی انتخاب و به بازۀ [\pnum{0}،\pnum{1}] نرمال می‌شوند. سپس ویژگی‌های مهندسی‌شده ساخته می‌شود و برای هر بیمار یا بردار جدولی \pnum{25}بعدی، یا دنبالۀ توکن مناسب رمزگذار \lr{ETHOS} با بردار نهفتۀ \pnum{128}بعدی تولید می‌گردد. برای شبیه‌سازی فدرال، کوهورت بر اساس تشخیص اصلی \lr{ICD-10} لایه‌بندی می‌شود.

\begin{table}[htbp]
\centering
\caption{گردش استخراج و پردازش دادۀ \lr{MIMIC\text{-}IV}}
\footnotesize
\setlength{\tabcolsep}{2pt}
\renewcommand{\arraystretch}{1.14}
\begin{tabularx}{\linewidth}{@{}R{1.8cm}Y Y@{}}
\toprule
مرحله & ورودی & خروجی \\
\midrule
استخراج & \texttt{\lr{patients}}، \texttt{\lr{admissions}}، \texttt{\lr{icustays}} & اقامت‌های \lr{ICU} با جمعیت‌شناسی \\
 & \texttt{\lr{labevents}}، \texttt{\lr{chartevents}} (از طریق \texttt{\lr{stay\_id}}) & آزمایشگاه و علائم حیاتی به‌ازای هر اقامت \\
انتخاب & \pnum{15} آزمایش + \pnum{7} علامت حیاتی & \lr{ItemID}ها از \texttt{\lr{d\_labitems}}، \texttt{\lr{d\_items}} \\
نرمال‌سازی & مقادیر خام & \lr{min-max} [\pnum{0}،\pnum{1}] \\
مهندسی & مقادیر نرمال‌شده & پراکسی \lr{SOFA}، شاخص شوک، شمار غیرطبیعی \\
بازنمایی & ویژگی‌های مهندسی‌شده & جدولی \pnum{25}-بعدی $\mid$ دنبالۀ توکن \\
رمزگذاری & دنبالۀ توکن & \lr{ETHOS} $\to$ بردار نهفتۀ \pnum{128}-بعدی \\
\midrule
\multicolumn{3}{@{}>{\raggedleft\arraybackslash}p{\linewidth}@{}}{\textbf{خروجی:} \pnum{2,000} بیمار؛ هر کدام \pnum{25} ویژگی به‌اضافه‌ی برچسب} \\
\bottomrule
\end{tabularx}
\label{tab:mimic_flow}
\end{table}

%==============================================================================
\section{جمعیت‌شناسی کوهورت}
%==============================================================================

جدول~\ref{tab:cohort_demographics_extended} آمار جمعیت‌شناختی کوهورت را نشان می‌دهد. در نمونۀ \pnum{2,000} بیمار، نرخ امکان‌پذیری درمان حدود \pnum{42}\% است و \pnum{58}\% نمونه‌ها در گروه ناممکن قرار می‌گیرند؛ این عدم‌تعادل نسبی نشان می‌دهد پیش‌بینی ترخیص \lr{ICU} مسئله‌ای ساده و کاملاً جداپذیر نیست. توزیع سن نیز به سمت سنین بالاتر کشیده شده و بیشتر بیماران در بازۀ \pnum{60} تا \pnum{80} سال قرار دارند.

\begin{table}[htbp]
\centering
\caption{جمعیت‌شناسی گستردۀ کوهورت}
\begin{tabular}{rr}
\toprule
معیار & مقدار \\
\midrule
کل بیماران & \pnum{2,000} \\
ممکن (برچسب=\pnum{1}) & \pnum{840} (\pnum{42}\%) \\
ناممکن (برچسب=\pnum{0}) & \pnum{1,160} (\pnum{58}\%) \\
میانگین سن (سال) & \pnum{64.2} \\
میانه سن (سال) & \pnum{67} \\
دامنۀ سن & \prange{18}{90} \\
زن & \pnum{44}\% \\
مرد & \pnum{56}\% \\
میانگین طول اقامت \lr{ICU} (ساعت) & \pnum{72.4} \\
میانه طول اقامت \lr{ICU} (ساعت) & \pnum{48} \\
\bottomrule
\end{tabular}
\label{tab:cohort_demographics_extended}
\end{table}

%==============================================================================
\section{توزیع گره بیماری}
%==============================================================================

جدول~\ref{tab:node_distribution} توزیع بیماران را در پنج گرۀ بیماری نشان می‌دهد. برای شفافیت و بازتولیدپذیری، جدول~\ref{tab:icd10_sample} نیز چند کد نمونه از \lr{ICD-10} را برای هر گره فهرست می‌کند.

\begin{table}[htbp]
\centering
\caption{نمونه‌هایی از کدهای \lr{ICD-10} به‌تفکیک گرۀ بیماری}
\footnotesize
\setlength{\tabcolsep}{2pt}
\renewcommand{\arraystretch}{1.14}
\begin{tabularx}{\linewidth}{@{}C{1cm}C{3.2cm}Y@{}}
\toprule
گره & کد & شرح \\
\midrule
\pnum{1} & \lr{A40}، \lr{A41} & استرپتوکوکی/سپسیس \\
\pnum{1} & \lr{R652} & سپسیس شدید \\
\pnum{2} & \lr{I21}، \lr{I22} & انفارکتوس حاد میوکارد \\
\pnum{2} & \lr{I50} & نارسایی قلبی \\
\pnum{3} & \lr{J18}، \lr{J44} & پنومونی، \lr{COPD} \\
\pnum{3} & \lr{J96} & نارسایی تنفسی \\
\pnum{4} & \lr{N17}، \lr{N18}، \lr{N19} & نارسایی حاد/مزمن کلیه \\
\pnum{5} & \lr{E10}، \lr{E11}، \lr{E13} & دیابت قندی \\
\bottomrule
\end{tabularx}
\label{tab:icd10_sample}
\end{table}
گره‌های سپسیس و قلبی از نظر تعداد نمونه بزرگ‌ترند، در حالی که گره‌های کلیوی و دیابت نمونه‌های کمتری دارند. همین تفاوت اندازه، استفاده از رویکردهای کم‌نمونه و فدرال را برای گروه‌های کم‌داده توجیه می‌کند.

\begin{table}[htbp]
\centering
\caption{توزیع بیماران به‌تفکیک گرۀ بیماری}
\footnotesize
\setlength{\tabcolsep}{2pt}
\renewcommand{\arraystretch}{1.14}
\begin{tabularx}{\linewidth}{@{}C{0.9cm}Y C{1.45cm}C{1.45cm}@{}}
\toprule
گره & گروه تشخیصی & بیماران & درصد ممکن \\
\midrule
\pnum{1} & سپسیس (\lr{A40}، \lr{A41}، \lr{R652}) & \pnum{412} & \pnum{38} \\
\pnum{2} & قلبی (\lr{I21}، \lr{I22}، \lr{I25}، \lr{I50}، و غیره) & \pnum{398} & \pnum{45} \\
\pnum{3} & تنفسی (\lr{J18}، \lr{J44}، \lr{J96}، و غیره) & \pnum{356} & \pnum{41} \\
\pnum{4} & کلیوی (\lr{N17}، \lr{N18}، \lr{N19}) & \pnum{418} & \pnum{39} \\
\pnum{5} & دیابت (\lr{E10}، \lr{E11}، \lr{E13}) & \pnum{416} & \pnum{46} \\
\bottomrule
\end{tabularx}
\label{tab:node_distribution}
\end{table}

%==============================================================================
\section{همبستگی آزمایشگاه و علائم حیاتی}
%==============================================================================

کراتینین با \lr{BUN} همبستگی بالایی دارد ($r > 0.7$) و هر دو شاخصی از عملکرد کلیه‌اند. هموگلوبین و پلاکت‌ها نیز بخشی از الگوی خون‌شناختی بیمار را منعکس می‌کنند، و در بیماران بحرانی گلوکز و لاکتات می‌توانند هم‌زمان افزایش یابند. این همبستگی‌ها برای مهندسی ویژگی مهم‌اند، زیرا نشان می‌دهند کدام ویژگی‌ها اطلاعات مکمل دارند و کدام‌ها ممکن است تا حدی افزونه باشند.

\begin{figure}[htbp]
\centering
\includegraphics[width=0.88\textwidth,keepaspectratio]{fig38_lab_correlation.png}
\caption{ماتریس همبستگی پیرسون برای آزمایشگاه‌ها و علائم حیاتی (ساختار گویای کوهورت آموزش).}
\label{fig:lab_correlation}
\end{figure}

%==============================================================================
\section{پراکسی \lr{SOFA} و شدت بالینی}
\label{sec:sofa_proxy}
%==============================================================================

پراکسی \lr{SOFA} از کراتینین، پلاکت‌ها، \lr{SpO$_2$}، و لاکتات محاسبه می‌شود. \lr{SOFA} بالاتر (اختلال اندام بدتر) با امکان‌پذیری پایین‌تر همبسته است، همان‌طور که بالینی انتظار دارد. میانگین پراکسی \lr{SOFA} برای بیماران ممکن \pnum{0.35} در برابر \pnum{0.52} برای ناممکن است ($p < 0.001$، رتبه‌جمع ویلکاکسون).

\begin{figure}[htbp]
\centering
\includegraphics[width=0.75\textwidth,keepaspectratio]{fig35_sofa_scatter.png}
\caption{پراکسی \lr{SOFA} در برابر کراتینین (نرمال‌شده) به‌تفکیک برچسب امکان‌پذیری؛ \lr{SOFA} و کراتینین بالاتر با ترخیص ناممکن همراه‌اند.}
\label{fig:sofa_scatter}
\end{figure}

%==============================================================================
\section{عملکرد لایه‌بندی‌شده بر اساس سن}
%==============================================================================

بیماران \pnum{65} سال و بالاتر نرخ امکان‌پذیری پایین‌تری دارند (\pnum{37}\%)؛ این نرخ برای گروه \pnum{18} تا \pnum{45} سال \pnum{51}\% و برای گروه \pnum{45} تا \pnum{65} سال \pnum{44}\% است. \lr{AUROC} مدل در گروه \pnum{65} سال و بالاتر نیز اندکی پایین‌تر از کل کوهورت است (\pnum{0.72} در برابر \pnum{0.76}). این تفاوت نشان می‌دهد که در صورت استقرار بالینی، ارزیابی انصاف و پایش عملکرد در زیرگروه‌های سنی ضروری است.

\begin{figure}[htbp]
\centering
\includegraphics[width=0.8\textwidth,keepaspectratio]{fig37_age_stratified.png}
\caption{\lr{AUROC} آزمون لایه‌بندی‌شده بر سن: \lr{Random Forest}، انباشت \lr{V4}، و خط پایۀ \lr{ETHOS} با یادگیری کم‌نمونه.}
\label{fig:age_stratified}
\end{figure}

%==============================================================================
\section{تحلیل خطا: کراتینین}
%==============================================================================

کراتینین از ویژگی‌های برتر است. طبقه‌بندی‌های اشتباه در بازه‌های مرزی (کراتینین نرمال‌شده \prange{0.3}{0.5}) خوشه‌بندی می‌شوند، جایی که عملکرد کلیه نه کاملاً طبیعی است و نه شدیداً مختل. در این بازه‌ها نرخ خطا حدود \pnum{40}\% در برابر \pnum{20}\% در بازه‌های شدید است.

\begin{figure}[htbp]
\centering
\includegraphics[width=0.78\textwidth,keepaspectratio]{fig36_error_creatinine.png}
\caption{نرخ خطای طبقه‌بندی به‌تفکیک بازۀ کراتینین (\lr{RF} در برابر \lr{ETHOS}) با شمار بیماران.}
\label{fig:error_creatinine}
\end{figure}

%==============================================================================
\section{تصویرسازی فضای نهفته (\lr{t-SNE})}
%==============================================================================

نمای \lr{t-SNE} از بردارهای نهفتۀ \lr{ETHOS} نشان می‌دهد بیماران ممکن و ناممکن تا حدی از هم جدا می‌شوند، اما هم‌پوشانی میان دو کلاس همچنان قابل توجه است. این الگو با عملکرد متوسط مدل کم‌نمونه سازگار است. امتیاز سیلوئت برای پارتیشن دوکلاسه برابر \pnum{0.12} است؛ یعنی ساختار خوشه‌ای ضعیف اما کاملاً تصادفی نیست.

\begin{figure}[htbp]
\centering
\includegraphics[width=0.75\textwidth,keepaspectratio]{fig27_tsne.png}
\caption{\lr{t-SNE} بردارهای نهفتۀ \lr{ETHOS} (بیماران آزمون): هم‌پوشانی جزئی کلاس‌های ممکن در برابر ناممکن.}
\label{fig:tsne_embeddings}
\end{figure}

%==============================================================================
\section{تحلیل نمونه}
%==============================================================================

در شبکه‌های نمونه‌ای، برای هر کلاس یک نمونه‌ی اولیه ساخته می‌شود که برابر میانگین بردارهای نهفتۀ مجموعۀ پشتیبان است. فاصله‌ی اقلیدسی میان دو نمونۀ اولیه در فضای نهفتۀ \pnum{128}بعدی برابر \pnum{3.42} است. این فاصله نشان می‌دهد سر کم‌نمونه تا حدی توان تفکیک دارد، اما جداسازی کلاس‌ها قوی نیست. بیماران پرس‌وجو بر اساس نزدیک‌ترین نمونه‌ی اولیه طبقه‌بندی می‌شوند.

\begin{figure}[htbp]
\centering
\includegraphics[width=0.82\textwidth,keepaspectratio]{fig28_prototype.png}
\caption{شماتیک مجموعه‌های پشتیبان، نمونه‌های کلاس، و طبقه‌بندی پرس‌وجو در فضای نهفته (پیش‌نمایش دوبعدی).}
\label{fig:prototype_illustration}
\end{figure}

%==============================================================================
\section{نمونۀ خط زمانی بیمار}
%==============================================================================

الگوهای زمانی، مانند کاهش کراتینین یا پایداری علائم حیاتی، برای برآورد امکان‌پذیری ترخیص اطلاعات مهمی فراهم می‌کنند. برای نمونه، بیماری که کراتینین او در \pnum{24} ساعت از \pnum{2.1} به \pnum{1.4} \lr{mg/dL} کاهش یافته و هم‌زمان ضربان قلب او در بازۀ \pnum{75} تا \pnum{85} \lr{bpm} و فشار خون سیستولیک او در بازۀ \pnum{120} تا \pnum{130} \lr{mmHg} پایدار مانده است، احتمال امکان‌پذیری بالاتری دارد (\pnum{0.78}). توکن‌سازی در این پژوهش چنین مسیرهای زمانی را به بازنمایی با طول ثابت تبدیل می‌کند.

\begin{figure}[htbp]
\centering
\includegraphics[width=0.9\textwidth,keepaspectratio]{fig32_patient_timeline.png}
\caption{نمونۀ مسیر \pnum{24} ساعت اول \lr{ICU}: ضربان قلب، فشار خون، و کراتینین (گویای).}
\label{fig:patient_timeline}
\end{figure}

%==============================================================================
\section{تعادل کلاس}
%==============================================================================

\lr{SMOTE} را برای مدیریت عدم‌تعادل \pnum{42}/\pnum{58} بررسی کردیم. در این وظیفه، وزن‌دهی کلاس در تابع زیان و تنظیم آستانه روی مجموعۀ اعتبارسنجی کافی بود و \lr{SMOTE} بهبود پایداری ایجاد نکرد. به‌نظر می‌رسد شدت عدم‌تعادل در این کوهورت متوسط است و برای بخش قابل توجهی از نمونه‌ها، مرز تصمیم با روش‌های ساده‌تر نیز قابل یادگیری است.

%==============================================================================
\section{فقدان داده و کیفیت داده}
%==============================================================================

در دادۀ بالینی، مقادیر گم‌شده پدیده‌ای رایج است؛ زیرا آزمایش‌ها با نظم یکسان برای همۀ بیماران درخواست نمی‌شوند و نیازهای بالینی هر بیمار متفاوت است. در این کوهورت، مقادیر آزمایشگاهی مانند لاکتات و تروپونین بیشتر از علائم حیاتی معمول مانند ضربان قلب و فشار خون گم می‌شوند. برای توکن‌سازی، در سری‌های زمانی از \lr{forward-fill} و برای ویژگی‌هایی که کاملاً گم شده‌اند از میانۀ جمعیت استفاده می‌شود. جدول~\ref{tab:missingness} نرخ فقدان را به‌تفکیک دسته‌ی ویژگی خلاصه می‌کند. ویژگی‌هایی با فقدان بیش از \pnum{30}\%، مانند لاکتات در برخی بیماران غیرسپتیک، در استقرار واقعی به پایش و رسیدگی دقیق‌تری نیاز دارند.

\begin{table}[htbp]
\centering
\caption{نرخ فقدان به‌تفکیک دستۀ ویژگی (\pnum{24} ساعت اول)}
\begin{tabular}{rcc}
\toprule
دستۀ ویژگی & درصد گم‌شده & جایگزینی \\
\midrule
آزمایشگاه‌های معمول (کراتینین، گلوکز) & \prange{5}{12} & میانه \\
آزمایشگاه‌های تخصصی (لاکتات، تروپونین) & \prange{18}{35} & میانه \\
علائم حیاتی (\lr{HR}، \lr{BP}، \lr{SpO2}) & \prange{2}{8} & \lr{Forward-fill} \\
مهندسی‌شده (\lr{SOFA}، شاخص شوک) & \pnum{0} & مشتق \\
\bottomrule
\end{tabular}
\label{tab:missingness}
\end{table}

%==============================================================================
\section{ارزیابی کیفیت داده}
%==============================================================================

کیفیت داده را با چند معیار سنجیدیم: (\pnum{1}) \textbf{کامل بودن}: سهم مقادیر غیرگم‌شده به‌ازای هر ویژگی؛ آزمایشگاه‌های معمول (کراتینین، سدیم) بیش از \pnum{90}\% کامل هستند؛ آزمایشگاه‌های تخصصی (لاکتات، تروپونین) \prange{65}{80}\%. (\pnum{2}) \textbf{باورپذیری}: مقادیر در بازه‌های بالینی (مثلاً کراتینین \prange{0.1}{30} \lr{mg/dL})؛ پرت‌ها را برش می‌دهیم. (\pnum{3}) \textbf{سازگاری}: بدون تناقض (مثلاً ترخیص پیش از پذیرش)؛ زمان‌بندی‌ها را اعتبارسنجی کردیم. (\pnum{4}) \textbf{نمایندگی}: کوهورت با توزیع سن و جنس \lr{MIMIC\text{-}IV} هم‌خوان است؛ نمونه‌برداری لایه‌ای ما تعادل بین گره‌های بیماری را تضمین می‌کند. این بررسی‌ها اعتبار نتایج تجربی را پشتیبانی می‌کنند.

%==============================================================================
\section{الگوهای زمانی}
%==============================================================================

روندهای زمانی در \pnum{24} ساعت نخست بستری بررسی شد. بیمارانی که آزمایش‌های آنان مسیر بهبود را نشان می‌دهد، برای مثال کاهش کراتینین یا پایداری گلوکز، معمولاً نرخ امکان‌پذیری بالاتری دارند. پراکسی \lr{SOFA} محاسبه‌شده از اوج مقادیر \pnum{24} ساعت نخست نیز با امکان‌پذیری همبسته است (بخش~\ref{sec:sofa_proxy}). در این پژوهش، اطلاعات زمانی به دنبالۀ توکنی با طول ثابت تبدیل می‌شود؛ در کار آینده می‌توان مدل‌سازی زمانی صریح، مانند \lr{LSTM} روی پنجره‌های ساعتی، را برای ثبت مستقیم‌تر روندها بررسی کرد.

%==============================================================================
\section{توزیع مقادیر آزمایشگاه و علائم حیاتی}
%==============================================================================

توزیع مقادیر کلیدی آزمایشگاه و علائم حیاتی در سراسر کوهورت بررسی شد. میانۀ کراتینین \pnum{1.1} \lr{mg/dL} است (\lr{IQR}: \pnum{0.8} تا \pnum{1.8}) و \pnum{12}\% بیماران مقداری بالاتر از \pnum{2.0} دارند که با اختلال عملکرد کلیه سازگار است. میانۀ گلوکز \pnum{140} \lr{mg/dL} است (\lr{IQR}: \pnum{110} تا \pnum{180}) و نوسان قند در بیماران بحرانی بالا دیده می‌شود. میانۀ هموگلوبین \pnum{9.2} \lr{g/dL} است (\lr{IQR}: \pnum{7.8} تا \pnum{10.5}) و \pnum{28}\% بیماران مقدار کمتر از \pnum{8.0} دارند. میانۀ لاکتات نیز \pnum{1.6} \lr{mmol/L} است (\lr{IQR}: \pnum{1.0} تا \pnum{2.5}) و لاکتات بالاتر از \pnum{2.0} در \pnum{35}\% بیماران با امکان‌پذیری پایین‌تر همراه است. این توزیع‌ها مبنای مهندسی ویژگی و تفسیر بالینی پیش‌بینی‌های مدل هستند. پراکسی \lr{SOFA} که از کراتینین، پلاکت‌ها، \lr{SpO$_2$} و لاکتات محاسبه می‌شود، میانۀ \pnum{0.42} دارد (\lr{IQR}: \pnum{0.28} تا \pnum{0.58})؛ مقادیر بالاتر نشان‌دهندۀ اختلال اندام بیشتر و نرخ امکان‌پذیری پایین‌تر است.

%==============================================================================
\section{نوع پذیرش و مقصد ترخیص}
%==============================================================================

امکان‌پذیری درمان بر اساس نوع پذیرش و مقصد ترخیص نیز تحلیل شد. پذیرش‌های اورژانسی نسبت به پذیرش‌های انتخابی نرخ امکان‌پذیری پایین‌تری دارند (\pnum{38}\% در برابر \pnum{48}\%)، که با شدت بالینی بیشتر سازگار است. در کل کوهورت، \pnum{22}\% بیماران به منزل، \pnum{18}\% به توان‌بخشی، \pnum{15}\% به مراقبت تخصصی پرستاری و \pnum{25}\% به مراقبت حاد دیگر ترخیص شده‌اند؛ مرگ بیمارستانی نیز \pnum{8}\% است. در تعریف برچسب این پژوهش، ترخیص به منزل، مراقبت خانگی و توان‌بخشی در گروه ممکن قرار می‌گیرند و سایر مقاصد در گروه ناممکن. این تعریف دودویی بخشی از ظرافت بالینی را از دست می‌دهد، اما مسئله را برای مدل‌سازی اولیه قابل مدیریت می‌کند.

%==============================================================================
\section{ساختار همبستگی و افزونگی ویژگی}
%==============================================================================

ساختار همبستگی، روابط بالینی مورد انتظار را نیز تأیید می‌کند. کراتینین و \lr{BUN} همبستگی قوی دارند (\mbox{\lr{$r > 0.7$}}) و هر دو عملکرد کلیه را بازتاب می‌دهند. هموگلوبین و هماتوکریت نیز همبستگی بسیار بالایی دارند (\mbox{\lr{$r > 0.9$}})، بنابراین از دید مدل می‌توانند اطلاعات افزونه ایجاد کنند. سدیم و کلر نیز با یکدیگر همبسته‌اند (\mbox{\lr{$r > 0.6$}}) و به تعادل الکترولیتی مربوط می‌شوند. این الگوها در انتخاب ویژگی اهمیت دارند؛ زیرا حذف ویژگی‌های بسیار افزونه می‌تواند مدل را ساده‌تر کند، بدون آنکه توان پیش‌بینی به‌طور محسوس کاهش یابد.

%==============================================================================
\section{جمع‌بندی فصل}
%==============================================================================

این فصل کوهورت \lr{MIMIC\text{-}IV} را از چند زاویه بررسی کرد: طرحوارۀ پایگاه داده و روند استخراج، جمعیت‌شناسی، توزیع گره‌های بیماری، همبستگی آزمایش‌ها و علائم حیاتی، پراکسی \lr{SOFA}، عملکرد به‌تفکیک سن، تحلیل خطا، بردارهای نهفته، تعادل کلاس، الگوهای فقدان، کیفیت داده، روندهای زمانی و ساختار همبستگی. مهم‌ترین یافته‌ها چنین‌اند: (\pnum{1}) کوهورت میان گره‌های بیماری نسبتاً متعادل است؛ (\pnum{2}) همبستگی‌های آزمایشگاهی با روابط بالینی شناخته‌شده هم‌خوانی دارد؛ (\pnum{3}) پراکسی \lr{SOFA} با امکان‌پذیری درمان رابطه دارد؛ (\pnum{4}) خطاهای طبقه‌بندی بیشتر در مقادیر مرزی آزمایشگاهی رخ می‌دهند؛ و (\pnum{5}) بردارهای نهفتۀ \lr{ETHOS} جداسازی جزئی اما نه کامل میان دو کلاس نشان می‌دهند. این نتایج مبنای انتخاب معماری، مهندسی ویژگی و تفسیر نتایج تجربی در فصل‌های بعد هستند.
